\usepackage[a4paper,inner=31mm,outer=27mm,top=28mm,bottom=31mm,marginparwidth=17mm]{geometry}
\usepackage{fontspec}
\defaultfontfeatures{Ligatures=TeX,Scale=MatchLowercase}
\setmainfont{Malgun Gothic}[
  ItalicFont={Malgun Gothic},
  ItalicFeatures={FakeSlant=0.2},
  BoldItalicFont={Malgun Gothic Bold},
  BoldItalicFeatures={FakeSlant=0.2}
]
\setsansfont{Malgun Gothic}
\setmonofont{TeX Gyre Cursor}
\XeTeXlinebreaklocale "ko"
\XeTeXlinebreakskip=0pt plus 1pt
\usepackage{microtype}
\usepackage{graphicx}
\usepackage{mathtools}
\usepackage{amsthm}
\usepackage{unicode-math}
\setmathfont{Latin Modern Math}
\usepackage{bookmark}
\usepackage{hyperref}
\hypersetup{hidelinks,unicode=true,pdfencoding=auto}

\setcounter{secnumdepth}{3}
\setcounter{tocdepth}{2}
\numberwithin{equation}{section}
\renewcommand{\contentsname}{목차}

\newcommand{\src}[1]{%
  \phantomsection\hypertarget{src:#1}{}%
  \marginpar{\raggedleft\footnotesize\sffamily #1}%
}
\newcommand{\sid}[1]{\phantomsection\hypertarget{fga:#1}{}}
\newcommand{\Osh}{\underline{O}}
\newcommand{\Zsh}{\underline{\mathbf Z}}
\newcommand{\Ccat}{\mathcal C}
\newcommand{\Homsh}{\underline{\operatorname{Hom}}}
\newcommand{\Extsh}{\underline{\operatorname{Ext}}}
\newcommand{\hsh}{\underline{h}}
\newcommand{\fquote}[1]{\textquotedbl#1\textquotedbl}
\newcommand{\seminar}[1]{%
  \begin{flushleft}\sffamily 부르바키 세미나\\#1\end{flushleft}%
}
\newcommand{\seminaren}[1]{%
  \begin{flushleft}\sffamily 부르바키 세미나\\#1\end{flushleft}%
}
\newcommand{\paperhead}[2]{%
  \begin{center}
    {\Large\bfseries #1\par}
    \vspace{0.6\baselineskip}
    #2
  \end{center}
  \vspace{0.8\baselineskip}
}

\theoremstyle{plain}
\newtheorem{proposition}{명제}[section]
\newtheorem{theoreme}[proposition]{정리}
\newtheorem{theorem}[proposition]{정리}
\newtheorem{corollaire}[proposition]{따름정리}
\newtheorem{corollary}[proposition]{따름정리}
\newtheorem{lemme}[proposition]{보조정리}
\newtheorem{lemma}[proposition]{보조정리}

\theoremstyle{definition}
\newtheorem{definition}[proposition]{정의}
\newtheorem{definitionen}[proposition]{정의}

\theoremstyle{remark}
\newtheorem*{remarque}{주}
\newtheorem*{remark}{주}

\renewcommand{\proofname}{증명}
\newcommand{\unit}[1]{%
  \chapter{#1}%
  \markboth{#1}{#1}%
}
\newcommand{\fsection}[2]{%
  \section*{#1. #2}%
  \addcontentsline{toc}{section}{#1. #2}%
}
\newcommand{\result}[2]{%
  \par\medskip\noindent\textbf{#1%
  \if\relax\detokenize{#2}\relax\else\space#2\fi. ---}\enspace
}
